\section{Le catalogue des trente-cinq}

\subsection{Durées et tailles}

\begin{figure}[H]
\centering
\includegraphics[width=\linewidth]{gold_trades_gantt}
\caption{Chaque trade, du premier au dernier. La longueur de la barre est la
durée de détention, son épaisseur le volume engagé, sa couleur le signe du
résultat. Les résultats supérieurs à 3\,000\,\$ en valeur absolue sont
annotés.}
\end{figure}

La lecture est frappante~: l'essentiel des barres est court et rouge, et
quelques barres très longues sont vertes. La durée médiane est de
\metric{5~jours}, la durée moyenne de \metric{34~jours} --- un écart qui dit à
lui seul que la distribution est dominée par sa queue.

\subsection{La cascade du résultat}

\begin{figure}[H]
\centering
\includegraphics[width=\linewidth]{gold_trades_waterfall}
\caption{Construction du résultat, dans l'ordre des clôtures. Chaque barre est
un trade~; la ligne sombre suit la balance de la sleeve.}
\end{figure}

Pendant les trois premières années, la balance oscille autour de son point de
départ~: le moteur ne gagne rien de significatif entre novembre~2021 et
fin~2024. Tout le résultat se construit en 2025 et au premier trimestre~2026,
sur deux positions. C'est une propriété du momentum, pas un accident~: la
stratégie attend une tendance et ne produit rien tant qu'il n'y en a pas.

\subsection{Durée contre résultat}

\begin{figure}[H]
\centering
\includegraphics[width=\linewidth]{gold_trades_scatter_duration}
\caption{Durée et résultat de chaque trade. La taille du point est le volume
engagé. L'échelle des durées est logarithmique.}
\end{figure}

La relation est nette. Les vingt-trois positions de moins de dix jours n'ont
jamais rapporté plus de \metric{529\,\$} l'une --- le vingt-et-un trades d'une
semaine ou moins totalisent $-$6\,188\,\$. À l'inverse, les cinq positions de
plus de cent jours totalisent $+$41\,096\,\$, et quatre d'entre elles portent
la quasi-totalité du résultat. La cinquième, ouverte en janvier~2022 et tenue
104~jours, perd 226\,\$~: la tenue longue n'est pas une garantie, c'est une
condition nécessaire.

\begin{insightbox}
Cette asymétrie est le mécanisme même du moteur, et elle a une conséquence
opérationnelle directe~: toute modification qui raccourcirait les positions ---
horizons plus courts, stop plus serré, sortie anticipée sur profit --- attaque
la seule partie de la distribution qui rapporte. C'est précisément ce qui avait
été observé lors du réglage des horizons~: la grille courte (15/30/60) semblait
meilleure côté recherche et s'effondrait à l'exécution.
\end{insightbox}

\subsection{La distribution}

\begin{figure}[H]
\centering
\includegraphics[width=\linewidth]{gold_trades_distribution}
\caption{Distribution des résultats par trade, gagnants et perdants séparés,
avec la dispersion en encart.}
\end{figure}

La forme est celle attendue d'un moteur de suivi de tendance~: un amas dense de
petites pertes près de zéro, et une queue droite très étirée. La moyenne
($+$918\,\$ par trade) se situe loin au-dessus de la médiane, qui est
\emph{négative}. Un lecteur qui jugerait ce moteur sur son trade médian
conclurait qu'il perd de l'argent.

% Généré par scripts/build_gold_trades_figures.py — NE PAS ÉDITER À LA MAIN.
\begin{table}[H]
\centering
\begin{tabular}{@{}lrrrr@{}}
\toprule
Année & Trades & Gagnants & Swap & Net \\
\midrule
2021 & 3 & 1 & -203 & -919 \\
2022 & 3 & 1 & -1,319 & -466 \\
2023 & 13 & 2 & -1,626 & -3,013 \\
2024 & 8 & 3 & -3,597 & +5,743 \\
2025 & 3 & 2 & -3,108 & +12,016 \\
2026 & 5 & 2 & -2,931 & +18,780 \\
\bottomrule
\end{tabular}
\caption{Résultat par année de clôture}
\end{table}

